\subsection{Analisi Within Country e Test dei Segni}\label{sec:methods_within}

%Per analizzare l'home advantage è stato effettuato un confronto within-country:
%la performance di ogni Paese che ha ospitato almeno una volta è stata
%confrontata con la media delle performance dello stesso Paese quando non ha
%ospitato. Obiettivo di tale analisi è verificare se esiste un vantaggio
%sistematico quando un Paese ospita le Olimpiadi.

%Calcolati i delta, ovvero le differenze fra la percentuale di medaglie vinte
%quando un Paese è ospitante e la media delle percentuali di medaglie vinte
%quando non lo è, è stato effettuato un Test dei Segni su tali delta. Il Test
%dei Segni è un test statistico non parametrico utilizzato per determinare se
%esiste una differenza sistematica tra coppie di osservazioni correlate. Tale
%test è statisticamente meno potente rispetto ad alternative come il t Test o il
%test di Wilcoxon. Tuttavia, questi due test richiedono che i dati provengano da
%distribuzioni normali o simmetriche, requisiti che abbiamo verificato che i
%nostri dati non rispettano.


%Per analizzare l'home advantage è stato effettuato un confronto within-country:
%la performance di ogni Paese che ha ospitato almeno una volta è stata
%confrontata con la media delle performance dello stesso Paese quando non ha
%ospitato.


% dettagli implementativi, non necessari

%Dal dataset originale sono quindi state estratte le righe relative a tutti i
%Paesi che hanno ospitato almeno una volta le Olimpiadi.

% unito sopra

%Fatto ciò, è stata calcolata la media delle percentuali di medaglie vinte da
%ogni Paese nelle edizioni in cui i paesi erano ospitanti e la media delle
%percentuali di medaglie vinte quando invece non lo erano. Sono stati poi
%calcolati i delta, ovvero le differenze fra queste due misure.


Sono state applicate due tecniche per constatare l'home advantage, confrontando,
per ogni Paese, la media delle medaglie (in percentuale) vinte in casa e
fuori.\\
Un'alisi within-country ha calcolato le differenze.\\



% IMPORTANTE: a cosa è stato applicato il sign test? Io ho scritto quello, ma
% non lo so io!

% non dobbiamo spiegare come funzionino i metodi, solo enunciarli e mostrarne i
% risultati e la validità.


%Il Test dei Segni è un test statistico non parametrico utilizzato per
%determinare se esiste una differenza sistematica tra coppie di osservazioni
%correlate. Viene spesso impiegato come alternativa al test t di Student o al
%test di Wilcoxon quando non è possibile garantire le assunzioni sulla
%distribuzione dei dati.
%
%In particolare, il test t di Student richiede che la distribuzione dei dati sia
%normale, mentre il test di Wilcoxon richiede che le differenze che si vogliono
%analizzare siano distribuite in maniera simmetrica.
%
%il Test dei Segni si basa esclusivamente sulla direzione della differenza tra
%le coppie. Per ogni coppia di dati (x, y), il test assegna:
%
%\begin{itemize} \item Un segno più se x $>$ y; \item Un segno meno se x $<$ y;
%    \item I casi di parità (differenza pari a zero) vengono solitamente esclusi
%    dall'analisi. \end{itemize}
%
%Sotto l'ipotesi nulla H0, ci si aspetta che il numero di segni positivi e
%negativi sia approssimativamente uguale, seguendo una distribuzione binomiale
%con probabilità p-value pari a 0.5.
%
%Sebbene sia molto flessibile, il Test dei Segni ha una minore potenza
%statistica rispetto al test di Wilcoxon, poiché ignora l'entità numerica delle
%differenze. Tuttavia, in presenza di distribuzioni fortemente non simmetriche,
%la sua affidabilità lo rende preferibile per evitare errori di tipo I (falsi
%positivi).
%
%Nel nostro caso, non conoscendo la distribuzione dei dati e avendo verificato
%una forte asimmetria, non è stato possibile usare il test t di Student o il
%test di Wilcoxon. Le differenze analizzate dal test dei segni sono i delta
%calcolati nel Within-Country Comparison.


% come è stata verificata l'asimmetria? Se non a occhio ma con un metodo, è
% IMPORTANTE dirlo!

Il Test dei Segni, poi, prendendo in input queste ultime, ritorna un segno
(positivo o negativo).\\
Le (più espressive) alternative Student's T-test e Wilcoxon (che esprimono anche
il valore della differenza) non potevano essere applicate, richiedendo,
rispettivamente, una distribuzione normale dei dati e una simmetrica delle
differenze.